\documentclass[11pt]{article}

\usepackage[margin=1in]{geometry}
\usepackage[T1]{fontenc}
\usepackage[utf8]{inputenc}
\usepackage{lmodern}
\usepackage{microtype}
\usepackage{amsmath,amssymb,amsthm}
\usepackage{booktabs}
\usepackage{enumitem}
\usepackage{xurl}
\usepackage{listings}
\usepackage{tabularx}
\usepackage[hidelinks]{hyperref}

\lstset{basicstyle=\footnotesize\ttfamily,breaklines=true,frame=single,columns=fullflexible}

\newtheorem{definition}{Definition}
\newtheorem{remark}{Remark}

\title{ABOM and OINL for Anonymity Claims:\\
Proof-Carrying Odds\,--\,Inflation Budgets Without Publishing Full Artifacts}
\author{}
\date{March 21, 2026 (archive v498)}

\begin{document}
\maketitle

\begin{abstract}
This series publishes brief, referee-grade papers, but a credible anonymity claim typically depends on a large artifact surface: measurement scripts, traces, compiler knobs, certificates, and verifier outputs.
Publicly shipping \emph{all} artifacts is often infeasible (space, privacy, operational constraints).

We propose two lightweight documentation objects that make claims checkable even in a source-only release.
An \emph{Anonymity Bill of Materials} (ABOM) is the signed claim object: it commits (by digest) to the claim semantics, artifact roots, and replay surfaces behind one privacy statement.
An \emph{Odds-Inflation Nutrition Label} (OINL) is the reader-facing diff surface extracted from that ABOM: it states the protected interface (ENF-ID), the per-epoch budget representation, and the intended window/threat declaration (TW-ID).
Together they let others diff claim objects across releases, verify signatures and digests, and request the full artifact archive only when needed, while leaving lineaged release receipts and publication policy to the separate release-property paper.
\end{abstract}

\paragraph{Scope.}
We isolate and formalize \emph{ABOM and OINL for Anonymity Claims} as a reusable contract surface in the anonymity / Anonymous-DHT series.

\paragraph{Imports.}
We treat odds inflation and endpoint sizing as imported (Anonymity~A; Synthesis~5).\cite{series:oddsinflation,series:endpointbridge}
We adopt the contract-object vocabulary (Synthesis~0) as the schema surface to which ABOM commits; CPPC manifests are a control-plane sibling of ABOM (AnonDHT~State~3).\cite{series:contractobjects,series:state3}
Interface declarations (exposure normal forms / ENF-ID) and threat/window declarations (TW-ID) are imported from the trace-interface and threat-window synthesis notes.\cite{series:traceifaces,series:threatwindows}
Replay-plan and state-declaration equivalence/drift rules are imported from Synthesis~18.\cite{series:planstateequiv}
For supply-chain style release receipts and transparency anchoring, see Release~A and W-Congestion-EQ.\cite{series:releaseproperty,series:wcongeq}

\paragraph{Artifact policy.} This archive is source-only; supporting artifacts (logs, traces, scripts, certificates, and verifier outputs) are available on request as a single bundle, committed by ABOM/OINL manifests (this paper) and governed by the series artifact policy (Synthesis~6).\cite{series:artifactpolicy}

\paragraph{Later terse citation posture.} Later terse papers should cite this note only when the live issue is the claim-object surface itself: which exact \texttt{claim\_id}, \texttt{exposure\_nf\_id}, \texttt{tw\_id}, \texttt{plan\_spec\_id}, \texttt{state\_decl\_id}, budget representation, and artifact-root digests were signed into one public claim packet. If the live issue is instead the mechanism-specific witness/accountant root, evaluator attempt logging and spending control, or cross-release lineage / receipt gating, stop at the imported mechanism owner, Eval~1, or Release~A rather than widening the claim-object layer into a second theorem or release dossier.\cite{series:evalnotary,series:releaseproperty}

\paragraph{Exports.}
Exports two schemas: (i) \emph{ABOM} (signed manifest of claim semantics and artifact digests) and
(ii) \emph{OINL} (a derived, human-readable label with claim id, projection, adjacency, budget representation, and horizon).
This paper deliberately stops before release-ledger lineage, transparency governance, and signed publication receipts; those are owned by Release~A.\cite{series:releaseproperty}

\section{Design goals}
ABOM borrows its ``bill of materials'' stance from software supply-chain practice: rather than publishing every artifact, we publish a signed manifest that commits to them and to their semantics.
This parallels SBOM standards (e.g., SPDX, CycloneDX) and attestation frameworks such as in-toto, but with anonymity-specific fields (threat windows, transcript projections, and budget accountants).\cite{ntiaSBOM2019,spdx301,cyclonedx,intoto2019}

ABOM/OINL are \emph{interfaces}, not proofs.
They aim to make three failure modes harder:
\begin{enumerate}[leftmargin=*]
\item \textbf{Silent drift:} knobs or evidence change without an auditable trail.
\item \textbf{Cherry-picked evidence:} a claim is supported by a subset of runs while other runs contradict it.
\item \textbf{Ambiguous semantics:} a budget is stated in a representation that downstream readers cannot interpret.
\end{enumerate}

\section{Public object split}
To keep the publication-facing papers non-redundant, we separate four objects that often get blurred together:
\begin{table}[t]
\centering
\small
\begin{tabular}{@{}p{2.3cm}p{4.7cm}p{4.9cm}@{}}
\toprule
\textbf{Object} & \textbf{What it owns} & \textbf{What it does not own}\\
\midrule
\textbf{ABOM} & The digest-bound claim object: semantics, artifact roots, replay/state hooks, and signer. & Cross-release comparison policy, receipt issuance, or transparency governance.\\
\addlinespace
\textbf{OINL} & The compact reader-facing summary used for diffing and quick review. & Proof objects, replay outputs, or release approval decisions.\\
\addlinespace
\textbf{Receipt / PVSA} & The signed release statement consumed by downstream tooling and dashboards. & Mechanism-specific witness construction or full artifact storage.\\
\addlinespace
\textbf{Artifact archive} & The heavy evidence bundle delivered on request. & Public-facing claim semantics unless bound by ABOM digests.\\
\bottomrule
\end{tabular}
\caption{Minimal ownership split for the publication-facing anonymity stack.}
\label{tab:object-split}
\end{table}
A useful mnemonic is: \emph{ABOM commits, OINL summarizes, receipts publish, archives substantiate}.

\paragraph{Series stack position.}
For the release-facing trio, Eval~1 owns the non-gameable attempt log and verifier bundle for one declared plan; ABOM/OINL owns the digest-bound claim object and the reader-facing diff surface extracted from it; Release~A owns cross-release lineage, signed receipts / PVSA, and fail-closed publication policy.\cite{series:evalnotary,series:releaseproperty}
The separation is intentional: a referee may audit whether one evaluation packet is real before deciding how that claim should be named, diffed, or published across releases.

\section{ABOM: the source of truth}
\begin{definition}[ABOM]\label{def:abom}
ABOM is a contract object in the shared vocabulary of Synthesis~0; this note records the anonymity-specific minimum surface.\cite{series:contractobjects}
An ABOM is a signed manifest that commits (by digest) to: a claim identifier and version; an ENF-ID naming the protected transcript projection; a TW-ID naming the threat/window declaration; an adjacency declaration; a budget representation; replay hooks (plan identifiers/spec digests); optional state-surface declarations when guarantees depend on long-lived control-plane state; artifact digests/availability flags; and a signature.
\end{definition}

\paragraph{Minimal schema sketch.}
The schema is intentionally permissive; it is a contract object for auditors.
A minimal ABOM (YAML/JSON) might contain:
\begin{lstlisting}
{
  "abom_version": "0.3",
  "claim_id": "anondht.destination_privacy.v1",
  "paper_ref": "series/anonymity_series/paperB_anonymous_dht_profiling",
  "exposure_nf_id": "sha256:<ENF-ID-hex>",
  "tw_id": "sha256:<TW-ID-hex>",
  "replay_hook": {"plan_id": "plan.replay.v1", "plan_spec_id": "sha256:<PLAN-SPEC-hex>"},
  "state_decl_id": "sha256:<STATE-DECL-hex>",
  "adjacency": "change-one-destination-class",
  "budget": {"type": "odds_inflation_bits", "epsilon_bits": 0.35, "delta": 1e-6, "unit": "lookup"},
  "artifacts": [
    {"name": "verifier_report.json", "sha256": "..."},
    {"name": "trace_bundle.tar.zst", "sha256": "...", "availability": "on_request"}
  ],
  "signing": {"scheme": "ed25519", "pubkey": "...", "signature": "..."}
}
\end{lstlisting}

\begin{remark}[Digests are the point]
ABOM allows a publication to be ``source-only'' while still being checkable: if an author later provides the private artifact archive, third parties can verify that the delivered artifacts match the published digests.
\end{remark}

\section{OINL: a label derived from ABOM}
OINL is a compact table extracted from the ABOM.
It is the object we expect most readers to consume and compare.

\begin{table}[t]
\centering
\small
\begin{tabularx}{0.97\linewidth}{@{}p{0.25\linewidth}X@{}}
\toprule
\textbf{OINL field} & \textbf{Meaning}\\
\midrule
Claim id & Stable string used for diffing and change control. \\
Exposure (ENF-ID) & \texttt{exposure\_nf\_id} naming the transcript interface whose budget is being claimed. \\
Adjacency & The privacy unit: what changes between worlds. \\
Budget type & Odds inflation, max-divergence, TV, or a declared tradeoff curve. \\
Per-epoch budget & $(\eta,\delta)$ or an equivalent declared representation. \\
Threat/window + usage & \texttt{tw\_id} plus the usage or horizon assumption and accounting regime (rulebook: Synthesis~19). \\
Owned proofs & Which papers or certificates provide the proof objects. \\
Artifact digests & Hash pointers to supporting evidence, retrievable on request. \\
\bottomrule
\end{tabularx}
\caption{OINL is a derived label: it summarizes (and links to) ABOM fields.}
\end{table}

\paragraph{A minimal public ABOM/OINL claim card.}
Table~\ref{tab:minimal-abom-card} is the smallest public packet later terse papers should need when the live question is simply \emph{what exact claim object was signed here?}

\begin{table}[t]
\centering
\small
\begin{tabularx}{0.97\linewidth}{@{}p{0.29\linewidth}X@{}}
\toprule
\textbf{Public row} & \textbf{Pinned value}\\
\midrule
Claim object & \texttt{claim\_id = anondht.lookup.destination\_privacy.v1}; OINL headline ``destination-class lookup privacy'' \\
Exposure + threat/window & \texttt{exposure\_nf\_id = sha256:5f3c\ldots b912}; \texttt{tw\_id = sha256:0ad8\ldots 2190} for the declared 30-day, up-to-\(10^4\)-lookups/day window \\
Replay / state hooks & \texttt{plan\_spec\_id = sha256:7cc1\ldots 01ab}; \texttt{state\_decl\_id = sha256:b0ae\ldots 90c1} \\
Adjacency + budget representation & \texttt{change-one-destination-class}; \texttt{odds\_inflation\_bits} with \((\eta,\delta)=(0.31,10^{-6})\) per lookup \\
Artifact roots & \texttt{verifier\_report.json = sha256:8a44\ldots 7f02}; \texttt{trace\_bundle.tar.zst = sha256:93c1\ldots 4c55}; availability \texttt{on\_request} \\
Signer & \texttt{ed25519} public key \texttt{K\_pub}; signed claim digest \texttt{sha256:4d77\ldots ee41} \\
Reader-facing OINL strip & ``destination-class lookup; \((\eta,\delta)=(0.31,10^{-6})\) per lookup; 30-day window; supporting bundle available on request'' \\
\bottomrule
\end{tabularx}
\caption{A minimal public ABOM/OINL claim card. This is the non-decorative public answer later terse papers can quote directly when the live issue is the signed claim object rather than evaluator spending, theorem-root arithmetic, or release lineage.}
\label{tab:minimal-abom-card}
\end{table}

Table~\ref{tab:minimal-abom-card} is deliberately narrow: it fixes the claim semantics, the replay/state hooks, and the artifact-root digests without pretending to settle whether the release is lineage-comparable or whether the evaluator packet was well spent. Those wider questions belong to Release~A and Eval~1.\cite{series:releaseproperty,series:evalnotary}

\paragraph{OINL and ``referee-grade brevity.''}
The papers in this archive can stay short because OINL externalizes the long tail of artifacts.
A referee can demand the archive; a casual reader can still interpret the claim.

\paragraph{What counts as an OINL-changing drift event.}
At minimum, a publication should mint a new ABOM/OINL object when any of the following changes: the claim identifier, the protected interface (ENF-ID), the threat/window declaration (TW-ID), the adjacency relation, the budget representation/unit, or the digest roots for replay/state hooks.
These are exactly the fields that tell a reader \emph{what experiment the number refers to}.
Whether two such objects remain comparable across releases is a release-ledger question owned by Release~A; this paper only insists that the changed claim surface be made explicit.\cite{series:releaseproperty}

\paragraph{One tiny drift drill.} Suppose Release~\(R\) and Release~\(R'\) share the same \texttt{claim\_id}, \texttt{exposure\_nf\_id}, \texttt{tw\_id}, \texttt{plan\_spec\_id}, \texttt{state\_decl\_id}, and budget type, but \texttt{verifier\_report.json} is reissued and the exported budget improves from \((0.31,10^{-6})\) to \((0.28,10^{-6})\). A later terse paper may then say ``same claim surface; refreshed evidence and recomputed budget'' without reopening the release ledger. If instead \texttt{tw\_id} or \texttt{plan\_spec\_id} changes, the terse paper should say ``new ABOM/OINL entry'' and leave same-lineage questions to Synthesis~18 or Release~A.\cite{series:planstateequiv,series:releaseproperty}

\paragraph{A minimal ABOM cutover rule.}
ABOM and OINL should fail closed at the claim-object boundary itself. If the reader-facing OINL strip is unchanged and only artifact-root digests or verifier outputs refresh, a later terse paper may say ``same OINL surface; refreshed ABOM digest/evidence.'' The moment any OINL-defining field changes---\texttt{claim\_id}, \texttt{exposure\_nf\_id}, \texttt{tw\_id}, adjacency, budget representation, \texttt{plan\_spec\_id}, or \texttt{state\_decl\_id}---the old public strip no longer licenses a quiet comparison: mint a new ABOM/OINL entry first, then let the neighboring owner decide whether the new entry is theorem-equivalent, accountant-equivalent, evaluator-equivalent, or only release-lineage comparable.

\paragraph{Minimal reopen ladder.}
Once the question is no longer ``what exact digest-bound claim object was signed here?'', reopen the first owner whose field changed. Reopen Anonymity~A when the endpoint semantics or generic odds-inflation conversion changes; Anonymity~B when the anonymous-DHT projection, secret scope, or local equalization witness changes; Boss Fight~A when the repeated-use window or usage accountant changes under the same endpoint semantics; State~1 when \texttt{state\_decl\_id} changes; Operational~A or Operational~B when retry discipline, runtime conformance, or replay-plan semantics change; Eval~1 when the live dispute is attempt logging, spending control, or verifier-bundle validity; and Release~A when the live issue is cross-release lineage, receipt / PVSA issuance, or publication gating. This note therefore exports one compact public stop rule: either the same OINL surface is being re-supported by a new ABOM digest, or a new claim-object entry must be minted before any higher-layer comparison is honest.\cite{series:bossfightA,series:profiling,series:state1,series:opretries,series:runtimeprivacy,series:evalnotary,series:releaseproperty}

\section{Interfacing with receipts, certificates, and transparency}

Receipt line items (Synthesis~12) are the publication-level companion of ABOM/OINL: each receipt line item repeats the same claim fields (\texttt{exposure\_nf\_id}, \texttt{tw\_id}, and when applicable \texttt{plan\_spec\_id} and \texttt{state\_decl\_id}) so that a reader can compare the OINL against the receipt without opening private artifacts.\cite{series:receiptitems}
ABOM/OINL therefore answer ``what exactly is this claim object?'', while Release~A answers ``how is that claim released, lineage-tagged, and policy-gated?''\cite{series:releaseproperty}

\paragraph{Horizon claims are structured.}
If an ABOM/OINL entry quotes a horizon-$n$ total ("per day", "for $Q$ lookups"), it should also name the usage assumption and bind the accountant via \texttt{plan\_spec\_id} (and \texttt{state\_decl\_id} when state enters), following the receipt-accounting rulebook.\cite{series:accountingrulebook}


\paragraph{A drift rule worth making explicit.}
If \texttt{plan\_spec\_id} or \texttt{state\_decl\_id} changes, the publication must mint a new receipt/ABOM entry that pins the new digest-bound semantics.
Whether the underlying privacy claim needs recertification depends on whether the change is merely ``publication-only'' (re-run the verifier) or claim-affecting; we centralize the taxonomy in Synthesis~18 so other papers do not restate it.\cite{series:planstateequiv}

ABOM is compatible with (and intentionally defers details to) the series' receipt/certificate pipeline:
\begin{itemize}[leftmargin=*]
\item \textbf{Release receipts:} Release~A defines a ledger bundle and a compact PVSA summary attestation.\cite{series:releaseproperty}
ABOM can point to the bundle and include the PVSA digest.
\item \textbf{Congestion certificates:} W-Congestion-EQ defines congestion-specific certificate fields and log anchoring patterns.\cite{series:wcongeq}
ABOM can embed the certificate digest and an inclusion-proof digest.
\item \textbf{Optimization certificates:} BossFight~C and the 2026-01-29 wiki post treat dual certificates as proof objects for TV privacy designs.
ABOM can include the primal/dual transcript digests.
\end{itemize}

\section{Takeaways}
\begin{itemize}[leftmargin=*]
\item \textbf{ABOM is the commitment:} a signed manifest that pins (by digest) the artifacts and the \emph{semantics} of the claim (projection, adjacency, budget representation, horizon), plus any replay and state surfaces the guarantee depends on.
\item \textbf{OINL is the diffable summary:} a compact label extracted from ABOM so reviewers and downstream users can compare claims across releases without unpacking the full artifact archive.
\item \textbf{Source-only can still be referee-grade:} publish ABOM/OINL in the paper; ship the heavy archive on request; let third parties verify the delivered bundle matches the published digests.
\end{itemize}

\section{Limitations and scope}
ABOM/OINL do not prevent lying: a publisher could sign a manifest for a claim that is false.
Their purpose is narrower: make claims \emph{auditable} and \emph{comparable} across releases, and make it straightforward for others to request and verify the full artifact archive.

\begin{thebibliography}{99}\small

\bibitem{series:artifactpolicy}
Anonymous.
\newblock Artifact and Disclosure Policy for the One-True-Archive (Source-Only Public Release).
\newblock Series synthesis note (Synthesis~6), 2026. (This archive.)

\bibitem{series:oddsinflation}
Anonymous.
\newblock Odds Inflation, Privacy Loss, and Endpoint Anonymity Bounds.
\newblock Anonymity series Paper~A, 2026. (This archive.)

\bibitem{series:endpointbridge}
Anonymous.
\newblock Endpoint Metrics Bridge: Odds Inflation, PML/MaxL, and Identification Sizing.
\newblock Series synthesis note (Synthesis~5), 2026. (This archive.)

\bibitem{series:evalnotary}
Anonymous.
\newblock Notarized Anonymity Certificates Under Retries.
\newblock Evaluation series Paper~1, 2026. (This archive.)

\bibitem{series:releaseproperty}
Anonymous.
\newblock Privacy as a Release Property: Proof-Carrying Anonymity Receipts for Anonymous DHTs.
\newblock Release and destination Paper~A, 2026. (This archive.)

\bibitem{series:wcongeq}
Anonymous.
\newblock W-Congestion-EQ: Proof-Carrying Congestion Budgets with Certificates and Transparency Logs for Anonymous Overlays.
\newblock Congestion series Paper~3, 2026. (This archive.)

\bibitem{series:traceifaces}
Anonymous.
\newblock Trace Interfaces for Anonymous DHTs: Observation Tiers, Committee-Contact Traces, and Exposure Normal Forms.
\newblock Series synthesis note (Synthesis~3), 2026. (This archive.)

\bibitem{series:threatwindows}
Anonymous.
\newblock Threat Models and Repetition Windows for Anonymous DHTs: A Short Declaration Checklist for Referee-Grade Budget Claims.
\newblock Series synthesis note (Synthesis~4), 2026. (This archive.)

\bibitem{series:accountingrulebook}
Anonymous.
\newblock Receipt Accounting and Horizon Composition for Anonymous-DHT Budgets.
\newblock Series synthesis note (Synthesis~19), 2026. (This archive.)

\bibitem{series:contractobjects}
Anonymous.
\newblock Contract Objects for Anonymous-DHT Privacy Budgets and Claims.
\newblock Series synthesis note (Synthesis~0), 2026. (This archive.)

\bibitem{series:receiptitems}
Anonymous.
\newblock Receipt Line Items for Anonymity Claims: A Minimal Tuple Schema for Published Budgets (Series Synthesis).
\newblock Series synthesis note (Synthesis~12), 2026. (This archive.)

\bibitem{series:planstateequiv}
Anonymous.
\newblock Digest-Bound Replay Plans and State Declarations: Equivalence, Minimality, and Change Control.
\newblock Series synthesis note (Synthesis~18), 2026. (This archive.)

\bibitem{series:bossfightA}
Anonymous.
\newblock The Boss Fight: A Receipt-Ready Accountant for Repeated Anonymity Use.
\newblock Bossfight series Paper~A, 2026. (This archive.)

\bibitem{series:profiling}
Anonymous.
\newblock Anonymous DHT Profiling and the Equalization Contract.
\newblock Anonymity series Paper~B, 2026. (This archive.)

\bibitem{series:opretries}
Anonymous.
\newblock Schedule-Only Retries for Anonymous DHTs.
\newblock Operational series Paper~A, 2026. (This archive.)

\bibitem{series:runtimeprivacy}
Anonymous.
\newblock Runtime Conformance and Auditable Trace Privacy.
\newblock Operational series Paper~B, 2026. (This archive.)

\bibitem{series:state1}
Anonymous.
\newblock State-Dependent Anonymity: Selection and State Evolution as an Attack Surface.
\newblock AnonDHT State series Paper~1, 2026. (This archive.)

\bibitem{series:state3}
Anonymous.
\newblock Closed-View Auditing for Stateful Anonymity: Odometers, Alarms, and Machine-Checkable CPPCs.
\newblock AnonDHT State series Paper~3, 2026. (This archive.)

\bibitem{ntiaSBOM2019}
National Telecommunications and Information Administration (NTIA).
\newblock Survey of Existing SBOM Formats and Standards.
\newblock Whitepaper, Oct.~2019.\;
\newblock \url{https://www.ntia.gov/sites/default/files/publications/ntia_sbom_formats_and_standards_whitepaper_-_version_20191025_0.pdf}.

\bibitem{cyclonedx}
OWASP Foundation.
\newblock CycloneDX Bill of Materials (BOM) Standard.
\newblock \url{https://cyclonedx.org/}. (Accessed Feb.~2026.)

\bibitem{spdx301}
The Linux Foundation (SPDX Project).
\newblock SPDX Specification v3.0.1.
\newblock 2024.\;\url{https://spdx.dev/wp-content/uploads/sites/31/2024/12/SPDX-3.0.1-1.pdf}.

\bibitem{intoto2019}
S.~Torres-Arias, H.~Afzali, T.~K.~Kuppusamy, R.~Curtmola, and J.~Cappos.
\newblock In-Toto: Providing farm-to-table guarantees for bits and bytes.
\newblock In \emph{USENIX Security}, 2019, pp.~1393--1410.\;
\newblock \url{https://www.usenix.org/conference/usenixsecurity19/presentation/torres-arias}.
\end{thebibliography}

\end{document}
